\begin{abstract}
Autonomous systems that hear can detect and locate events that cameras and lidar miss, such as a siren behind a building or a voice under rubble, yet audition remains the least developed sense on robots, drones and vehicles.
This review surveys acoustic perception for these three platform families across six tasks: localization and tracking, event detection, separation and enhancement, joint localization and detection, acoustic mapping, and audio-driven embodied behavior, tracing each from array signal processing to learned, multimodal and language-grounded models.
We argue that one problem sets the field apart from audio on static devices: \emph{ego-noise}, the sound a platform makes itself.
When it dominates, methods that model it explicitly, especially with the platform's own motor or joint state, outperform those that treat it as generic interference.
We review datasets, simulators, toolkits and evaluation protocols, and report reproducible case studies on three public datasets.
On a drone, a learned rotor-noise model raises in-flight localization accuracy within $10^\circ$ from 18\% to 70\%, motor speed helps most when the target is quiet, and the model gives about 10\,dB of beamforming gain.
Across five drones, the tonal share of ego-noise ranges from 1\% to 70\%. The highest values come from drones that were not flying, and rotor-speed variation in flight lowers the tonal share. A learned spatial model removes 6 to 13\,dB on every drone.
On a robot-head array without ego-noise, the same model lowers accuracy for moving talkers. This marks where such models stop helping.
We close with open problems, from hearing at very low signal-to-noise ratios and vibration-referenced noise cancellation to transfer across platforms and privacy.
\end{abstract}
